\documentclass{article}
\usepackage[utf8]{inputenc}
\usepackage{geometry}
\geometry{a4paper, margin=1in}
\usepackage{hyperref}
\usepackage{listings}
\usepackage{xcolor}
\usepackage{graphicx}
\usepackage{booktabs}
\usepackage{array}

\definecolor{codegreen}{rgb}{0,0.6,0}
\definecolor{codegray}{rgb}{0.5,0.5,0.5}
\definecolor{codepurple}{rgb}{0.58,0,0.82}
\definecolor{backcolour}{rgb}{0.95,0.95,0.92}

\lstdefinestyle{mystyle}{
    backgroundcolor=\color{backcolour},
    commentstyle=\color{codegreen},
    keywordstyle=\color{magenta},
    numberstyle=\tiny\color{codegray},
    stringstyle=\color{codepurple},
    basicstyle=\ttfamily\footnotesize,
    breakatwhitespace=false,
    breaklines=true,
    captionpos=b,
    keepspaces=true,
    numbers=left,
    numbersep=5pt,
    showspaces=false,
    showstringspaces=false,
    showtabs=false,
    tabsize=2
}

\lstset{style=mystyle}

\title{Comprehensive Python Learning Guide: From Beginner to Advanced}
\author{Kay Böhmer}
\date{\today}

\begin{document}

\maketitle
\tableofcontents
\newpage

\section{Comprehensive Python Learning Guide: From Beginner to Advanced}

- Version: 1.0 ( Free Edition )
- Author: [ Kay Böhmer ]
- Date: May 2026
- License: Creative Commons ( CC0-1.0 license )

\subsection{Table of Contents}

\begin{itemize}
    \item [Introduction to Python](\#introduction-to-python)
    \item [Python Basics](\#python-basics)
    \item [Control Flow](\#control-flow)
    \item [Functions and Modules](\#functions--modules)
    \item [Data Structures](\#data-structures)
    \item [File Handling](\#file-handling)
    \item [Object-Oriented Programming (OOP)](\#object-oriented-programming-oop)
    \item [Advanced Python Concepts](\#advanced-python-concepts)
    \item [Debugging and Testing](\#debugging--testing)
    \item [Working with Libraries and Frameworks](\#working-with-libraries--frameworks)
    \item [Practical Exercises and Projects](\#practical-exercises--projects)
    \item [Quizzes and Knowledge Checks](\#quizzes--knowledge-checks)
    \item [Interactive Learning Modules](\#interactive-learning-modules)
    \item [Common Mistakes and Best Practices](\#common-mistakes--best-practices)
    \item [Free Learning Resources](\#free-learning-resources)
\end{itemize}

\section{Introduction to Python}

Python is a high-level, interpreted, and general-purpose programming language known for its readability, versatility, and strong community support. It is widely used in industries like web development, data science, artificial intelligence, and automation.

\subsection{Why Learn Python?}

\begin{itemize}
    \item Readability: Clean and easy-to-understand syntax.
    \item Versatility: Used in web development, data science, AI, automation, and more.
    \item Community Support: Large ecosystem of libraries and frameworks.
    \item Beginner-Friendliness: Ideal for beginners due to its simplicity.
\end{itemize}

\subsection{Python in Industry \& Academia}

\begin{itemize}
    \item Industry: Used by companies like Google, Netflix, NASA, and Instagram.
    \item Academia: Taught in universities worldwide for introductory programming courses.
\end{itemize}

\subsection{Setting Up Python}

\subsubsection{Installing Python}

\begin{itemize}
    \item Windows: Download the installer from [python.org](https://www.python.org/downloads/) and run it.
    \item Linux: Use the package manager (e.g., sudo apt install python3 for Ubuntu).
    \item MacOS: Use Homebrew (brew install python) or download from [python.org](https://www.python.org/downloads/macos/).
\end{itemize}

\subsubsection{Choosing an IDE}

\begin{itemize}
    \item Beginner-Friendly: Thonny, IDLE, VS Code (with Python extension).
    \item Advanced: PyCharm, Jupyter Notebook, Spyder.
\end{itemize}

\subsubsection{Using Anaconda (Optional)}

Anaconda is a data science-focused distribution of Python that includes the Python interpreter, Jupyter Notebook, and popular libraries like NumPy, Pandas, and Matplotlib. Download from [anaconda.com](https://www.anaconda.com/products/distribution).

\section{Python Basics}

\subsection{Syntax \& Indentation}

Python uses indentation (whitespace) to define code blocks (instead of braces {}).

\begin{lstlisting}[language=Python]
if 5 > 2:
    print("Five is greater than two!")
\end{lstlisting}

\subsection{Variables \& Data Types}

Variables in Python are dynamically typed. Common data types include:

\begin{itemize}
    \item int: Integer (e.g., x = 10).
    \item float: Floating-point number (e.g., y = 3.14).
    \item str: String (e.g., name = "Alice").
    \item bool: Boolean (e.g., is_valid = True).
    \item list: Ordered, mutable collection (e.g., nums = [1, 2, 3]).
    \item tuple: Ordered, immutable collection (e.g., point = (1, 2)).
    \item dict: Key-value pairs (e.g., person = {"name": "Alice"}).
    \item set: Unordered, unique elements (e.g., unique_nums = {1, 2, 3}).
\end{itemize}

\begin{lstlisting}[language=Python]
# Variable assignment
age = 25                # int
price = 19.99           # float
name = "Kay"            # str
is_student = True       # bool
fruits = ["apple", "banana"]  # list
coordinates = (10, 20) # tuple
person = {"name": "Kay", "age": 25}  # dict
unique_numbers = {1, 2, 3}    # set
\end{lstlisting}

\subsection{Basic Operators}

\begin{table}[h]
\centering
\begin{tabular}{|l|l|l|}
\hline
\textbf{Operator} & \textbf{Example} & \textbf{Description} \\
\hline
+ & x + y & Addition \\
- & x - y & Subtraction \\
* & x * y & Multiplication \\
/ & x / y & Division (returns float) \\
\% & x \% y & Modulus (remainder) \\
** & x ** y & Exponentiation \\
// & x // y & Floor division \\
== & x == y & Equal to \\
!= & x != y & Not equal to \\
> & x > y & Greater than \\
< & x < y & Less than \\
and & x and y & Logical AND \\
or & x or y & Logical OR \\
not & not x & Logical NOT \\
\hline
\end{tabular}
\end{table}

\subsection{Input/Output}

\begin{lstlisting}[language=Python]
name = input("Enter your name: ")
print(f"Hello, {name}!")
\end{lstlisting}

\section{Control Flow}

\subsection{Conditional Statements}

\begin{lstlisting}[language=Python]
age = 18
if age < 13:
    print("Child")
elif age < 18:
    print("Teenager")
else:
    print("Adult")
\end{lstlisting}

\subsection{Loops}

\subsubsection{for Loop}

\begin{lstlisting}[language=Python]
fruits = ["apple", "banana", "cherry"]
for fruit in fruits:
    print(fruit)
\end{lstlisting}

\subsubsection{while Loop}

\begin{lstlisting}[language=Python]
count = 0
while count < 5:
    print(count)
    count += 1
\end{lstlisting}

\section{Functions \& Modules}

\subsection{Defining Functions}

\begin{lstlisting}[language=Python]
def greet(name):
    """This function greets the person passed in as a parameter"""
    return f"Hello, {name}!"

print(greet("Kay"))  # Output: Hello, Kay!
\end{lstlisting}

\subsection{Importing Modules}

\begin{lstlisting}[language=Python]
import math
print(math.sqrt(16))  # Output: 4.0

from math import pi
print(pi)  # Output: 3.141592653589793
\end{lstlisting}

\section{Data Structures}

\subsection{Lists}

\begin{lstlisting}[language=Python]
fruits = ["apple", "banana", "cherry"]
fruits.append("orange")
print(fruits)  # Output: ['apple', 'banana', 'cherry', 'orange']
\end{lstlisting}

\subsection{Dictionaries}

\begin{lstlisting}[language=Python]
person = {"name": "Kay", "age": 25}
print(person["name"])  # Output: Kay
\end{lstlisting}

\section{OOP in Python}

\subsection{Classes \& Objects}

\begin{lstlisting}[language=Python]
class Person:
    def __init__(self, name, age):
        self.name = name
        self.age = age

    def greet(self):
        return f"Hello, I'm {self.name}!"

person = Person("Kay", 25)
print(person.greet())  # Output: Hello, I'm Kay!
\end{lstlisting}

\section{Advanced Python Concepts}

\subsection{Decorators}

\begin{lstlisting}[language=Python]
def logger(func):
    def wrapper(*args, **kwargs):
        print(f"Calling {func.__name__}")
        return func(*args, **kwargs)
    return wrapper

@logger
def add(a, b):
    return a + b

print(add(2, 3))  # Output: Calling add \n 5
\end{lstlisting}

\section{Practical Exercises \& Projects}

\subsection{Beginner Projects}

\begin{itemize}
    \item Hello World
    \item Temperature Converter
    \item Simple Calculator
\end{itemize}

\subsection{Intermediate Projects}

\begin{itemize}
    \item Web Scraper
    \item File Encryption/Decryption
\end{itemize}

\section{Quizzes \& Knowledge Checks}

\subsection{Python Basics Quiz}

\begin{enumerate}
    \item What is the output of \texttt{print(10 \% 3)}?
    \begin{itemize}
        \item A) 1
        \item B) 2
        \item C) 3
        \item D) 0
    \end{itemize}
    \textbf{Answer:} A) 1
\end{enumerate}

\section{Free Learning Resources}

\begin{table}[h]
\centering
\begin{tabular}{|l|l|l|}
\hline
\textbf{Resource} & \textbf{Description} & \textbf{Link} \\
\hline
Official Python Docs & Comprehensive tutorials \& references & \href{https://docs.python.org/3/}{python.org} \\
W3Schools Python & Beginner-friendly examples & \href{https://www.w3schools.com/python/}{w3schools.com/python} \\
Codecademy & Interactive Python course & \href{https://www.codecademy.com/learn/learn-python-3}{codecademy.com} \\
LearnPython.org & Free interactive tutorial & \href{https://www.learnpython.org/}{learnpython.org} \\
\hline
\end{tabular}
\end{table}

\section{Appendices}

\subsection{Python Cheat Sheet}

\begin{lstlisting}[language=Python]
# Variables
x = 10
name = "Kay"

# Lists
lst = [1, 2, 3]
lst.append(4)

# Functions
def add(a, b):
    return a + b
\end{lstlisting}

\end{document}
